\documentclass{article}
%packages
\usepackage{fullpage}
\usepackage{xspace}
\usepackage{graphicx,color}
\usepackage{hyperref}
\usepackage{graphicx}
\usepackage[dvipsnames]{xcolor}
\usepackage{amsmath, amsthm, amscd, amsfonts, amssymb}
\usepackage{graphicx}
\usepackage{stmaryrd}
\usepackage{quiver}
\usepackage{float}
\usepackage{mathtools}
\usepackage{mathrsfs}
\usepackage{multicol}
\usepackage{tikz-cd}
\usepackage[shortlabels]{enumitem}
\usepackage{stackrel}
\usepackage{cancel}
\usepackage[framemethod = TikZ]{mdframed}

%theorems
\newtheorem{theorem}{Teorema}[section]

\newtheorem{algorithm}{Algorithm}
\newtheorem{axiom}{Axioma}
\newtheorem{case}{Case}
\newtheorem{claim}{Afirmación}
\newtheorem{conclusion}{Conclusión}
\newtheorem{condition}{Condition}
\newtheorem{conjecture}[theorem]{Conjetura}
\newtheorem{corollary}[theorem]{Corolario}
\newtheorem{criterion}{Criterion}
\newtheorem{proposition}[theorem]{Proposición}
\newtheorem{lemma}[theorem]{Lema}

\theoremstyle{definition}
\newtheorem{definition}[theorem]{Definición}

\newtheorem{exercise}{Ejercicio}

\newtheorem*{notation}{Notación}
\newtheorem{problem}{Problema}[section]
\newtheorem*{obs}{Observación}
\newtheorem*{addendum}{Addendum}

\newtheorem*{tarea}{\wea{Tarea}}

\theoremstyle{remark}
\newtheorem{remark}[theorem]{Remark}
\newtheorem*{solution}{\textbf{Solución}}
\newtheorem{summary}{Summary}
\newtheorem{example}[theorem]{\textbf{Ejemplo}}

\numberwithin{equation}{section}

\renewcommand*{\proofname}{\textbf{Demostración}}

\surroundwithmdframed[
  topline=false,
  rightline=false,
  bottomline=false,
  leftmargin=\parindent,
  skipabove=\medskipamount,
  skipbelow=\medskipamount
]{proof}

\surroundwithmdframed[
  topline=false,
  rightline=false,
  bottomline=false,
  leftmargin=\parindent,
  skipabove=\medskipamount,
  skipbelow=\medskipamount
]{solution}

%commands
\newcommand{\A}{\mathbb{A}}
\newcommand{\B}{\mathcal{B}}
\newcommand{\C}{\mathbb{C}}
\newcommand{\D}{\mathbb{D}}
\newcommand{\F}{\mathbb{F}}
\newcommand{\N}{\mathbb{N}}
\renewcommand{\O}{\mathcal{O}}
\renewcommand{\P}{\mathcal{P}}
\newcommand{\Q}{\mathbb{Q}}
\newcommand{\R}{\mathbb{R}}
\renewcommand{\S}{\mathbb{S}}
\newcommand{\U}{\mathcal{U}}
\newcommand{\Z}{\mathbb{Z}}

\newcommand{\Acal}{\mathcal{A}}
\newcommand{\Bcal}{\mathcal{B}}
\newcommand{\Ccal}{\mathcal{C}}
\newcommand{\Dcal}{\mathcal{D}}
\newcommand{\Ecal}{\mathcal{E}}
\newcommand{\Fcal}{\mathcal{F}}
\newcommand{\Gcal}{\mathcal{G}}
\newcommand{\Hcal}{\mathcal{H}}
\newcommand{\Ical}{\mathcal{I}}
\newcommand{\Jcal}{\mathcal{J}}
\newcommand{\Kcal}{\mathcal{K}}
\newcommand{\Lcal}{\mathcal{L}}
\newcommand{\Mcal}{\mathcal{M}}
\newcommand{\Ncal}{\mathcal{N}}
\newcommand{\Qcal}{\mathcal{Q}}
\newcommand{\Rcal}{\mathcal{R}}
\newcommand{\Scal}{\mathcal{S}}
\newcommand{\Tcal}{\mathcal{T}}
\newcommand{\Ucal}{\mathcal{U}}
\newcommand{\Vcal}{\mathcal{V}}
\newcommand{\Wcal}{\mathcal{W}}
\newcommand{\Xcal}{\mathcal{X}}
\newcommand{\Ycal}{\mathcal{Y}}
\newcommand{\Zcal}{\mathcal{Z}}

\renewcommand{\a}{\alpha}
\renewcommand{\b}{\beta}
\renewcommand{\d}{\delta}
\newcommand{\e}{\varepsilon}
\renewcommand{\t}{\tau}

\newcommand{\vphi}{\varphi}
\newcommand{\oo}{\infty}
\newcommand{\isom}{\cong}
\newcommand{\longto}{\longrightarrow}
\newcommand{\embbed}{\hookrightarrow}
\renewcommand{\-}{\setminus}

\renewcommand{\bar}{\overline}
\newcommand{\sgn}{\mathrm{sgn}}
\newcommand{\id}{\mathrm{id}}
\newcommand{\im}{\mathrm{im}\,}
\newcommand{\dist}{\mathrm{dist}}

\newcommand{\normal}{\unlhd}

\renewcommand{\tilde}{\widetilde}
\renewcommand{\hat}{\widehat}

\newcommand{\abs}[1]{\left| #1 \right|}
\newcommand{\norm}[1]{\left\lVert #1 \right\rVert}
\newcommand{\br}[1]{\left( #1 \right)}
\newcommand{\set}[1]{\left\{ #1 \right\}}
\newcommand{\cor}[1]{\left[ #1 \right]}
\newcommand{\gen}[1]{\left\langle #1 \right\rangle}
\newcommand{\matr}[1]{\begin{pmatrix}#1\end{pmatrix}}

\newcommand{\znz}[1]{\frac{\Z}{#1 \Z}}
\newcommand{\zmz}[1]{\Z/ #1 \Z}

% Por si me equivoco al escribir
\newcommand{\rac}[2]{\frac{#1}{#2}}
\newcommand{\olon}{\colon}
\newcommand{\ubseteq}{\subseteq}
\newcommand{\ubset}{\subset}
\newcommand{\irc}{\circ}

\newcommand{\ds}{\displaystyle}

\renewcommand*\contentsname{Contenidos}

\newcommand{\wea}[1]{\textcolor{BlueViolet}{\textbf{\emph{#1}}}}
\newcommand{\fecha}[1]{\begin{flushright} \underline{#1} \end{flushright}}

\newcommand{\dem}[1]{\textbf{Demostración (#1)}}

%title
\title{MPG3200 -- Análisis I}
\author{Felipe Inostroza \\ Prof. Mariel Sáez}
\date{Primer Semestre 2026}

\begin{document}

\maketitle
\tableofcontents

\newpage
\fecha{Clase 1. Lunes 9 de Marzo}
\section{Topología}
\subsection{Topologías y espacios métricos}

\begin{definition}
  Dado un conjunto $X$, un conjunto
  \[\tau \subseteq \P(X) = \{\O \colon \O \subseteq X\}\]
  se dice \wea{topología en $X$} si
  \begin{enumerate}
    \item $\varnothing, X \in \tau$;
    \item Si $\O_1, \O_2 \in \tau \implies \O_1 \cap \O_2 \in \tau$;
    \item $\{\O_\a\}_{\a\in I} \subseteq \tau \implies \bigcup_{\a \in I} \O_\a \in \tau$.
  \end{enumerate}
\end{definition}

\begin{notation}
  \begin{itemize}
    \item Si $\O \in \t$, $\O$ se llama \wea{abierto}.
    \item $(X,\tau)$ es un \wea{espacio topológico}.
    \item Si $a \in X$ y $\O \in \t$ tal que $a \in \O$, $\O$ se llama \wea{vecindad de $a$}.
  \end{itemize}
\end{notation}

\begin{example}
  \begin{enumerate}
    \item Para cualquier conjunto $X$ se tiene que $\t = \{\varnothing,X\}$ es una topología de $X$ (\wea{topología trivial}).
    \item $\t = \P(X)$ es la \wea{topología discreta}.
    \item Si $(X,d)$ es un espacio métrico, una topología es
    \[\t = \{\O \subseteq X \colon \forall a \in \O, \exists r > 0 \colon B_r(A) \subseteq \O\}.\]
    Esto incluye $\R^n$, espacios normados, espacios de Hilbert, etc.
  \end{enumerate}
\end{example}

\begin{proposition}
  $E \ubseteq X$ es abierto si y solo si $\forall a \in E$, $\exists \O_a$ vecindad de $a$ tal que $\O_a \subseteq E$.
\end{proposition}
\begin{proof}
  \begin{itemize}
    \item \boxed\implies
    Si $a \in E$ y $E$ es abierto, entonces $E$ es vecindad de $a$ y $E \subseteq E$. \checkmark
    \item \boxed\impliedby
    Por demostrar: $E$ abierto.
    Si $a \in E$, por hipótesis $\exists \O_a$ tal que $a \in \O_a \subseteq E$, y por lo tanto
    \[\bigcup_{a \in E} \O_a \subseteq E \subseteq \bigcup_{a \in E} \O_a.\]
    Por lo tanto $E = \bigcup_{a \in E} \O_a$, y al ser unión de abiertos, tenemos que $E$ es abierto.
  \end{itemize}
\end{proof}

\begin{definition}
  La \wea{clausura} de un conjunto $E$ está dada por
  \[\bar{E} = \{a \in X \colon \forall \O_a \ni a, \ \O_a \cap E \neq \varnothing\}.\]
\end{definition}

\begin{definition}
  $E$ es cerrado si y solo si $\bar{E} = E$.
\end{definition}

\begin{proposition}
  $E$ es cerrado si y solo si $X \setminus E$ es abierto.
\end{proposition}
\begin{proof}
  \
  \begin{itemize}
    \item \boxed{\implies} Sea $a \in X \setminus E$. Si $\forall \O_a$ vecindad de $a$
    \begin{align*}
      \O_a \not\subseteq X \setminus E
      &\iff \exists b \in \O_a \cap E \\
      &\iff \O_a \cap E \neq \varnothing \\
      &\iff a \in \bar E = E.
    \end{align*}
    Pero esto contradice que $a \in X \setminus E$. Entonces existe $\O_a$ tal que $\O_a \subseteq X \setminus E$, por lo que $X \setminus E$ es abierto.

    \item \boxed{\impliedby} Por demostrar: $\bar E = E$.
    Si $a \in E$, entonces $\forall \O_a$ vecindad de $a$ se tiene que $$a \in \O_a \cap E$$ y entonces $E \subseteq \bar E$. Veamos ahora que $\bar E \subseteq E$.

    Sea $a \in \bar E$, entonces para toda vecindad $\O_a$ se tiene que
    \[\O_a \cap E \neq \varnothing. \tag{*}\]
    Supongamos que $a \in \bar E \setminus E$, entonces $a \in X \setminus E$. Como $X \setminus E$ es abierto, existe $\O_*$ vecindad de $a$ tal que $\O_* \subseteq X \setminus E$. Entonces $\O_* \cap E = \varnothing$, lo que contradice $(*)$.
    Concluimos que $\bar E \setminus E = \varnothing$.
  \end{itemize}
\end{proof}

\begin{definition}
  [Topología inducida]
  Si $(X,\t)$ es un espacio topológico e $Y \subseteq X$, la \wea{topología inducida en $Y$} es
  \[\t|_Y = \{\O \cap Y \colon \O \in \t\}.\]
\end{definition}

\begin{exercise}
  Verificar que $\t|_Y$ es topología.
\end{exercise}

\begin{obs}
  Con la topología inducida, los abiertos de $Y$ no son necesariamente abiertos en $X$.
\end{obs}

\begin{example}
  En $\R$ con la topología usual, si tomamos $Y = [0,2]$ entonces $[0,1)$ es abierto, ya que $[0,1) = [0,2] \cap (-1,1)$ y $(-1,1)$ es abierto en $\R$.
\end{example}

\begin{example}
  Consideremos la \wea{2--esfera}
  \[\S^2 = \{(x,y,z) \in \R^3 \colon x^+ y^2 + z^= 1\}.\]
  El conjunto
  \[\{(x,y,z) \in \S^2 \colon z>0\}\]
  es abierto en la topología inducida por $\R^3$.
\end{example}

\begin{definition}
  Sea $P \in X$. Decimos que $\B_p \subseteq \t$ es una \wea{base de la topología $\t$ en $p$} si
  \begin{enumerate}[(i)]
    \item Todos los elementos de $\B_p$ son vecindades de $p$;
    \item Para toda vecindad $U \ni p$, existe $B \in \B_p$ tal que $B \subseteq U$.
  \end{enumerate}
  Decimos que $\B$ es \wea{base de la topología} si $\forall p \in X$, $\exists \B_p \subseteq \B$ base de la topología en $p$.
\end{definition}

\begin{example}
  Si $(X,d)$ es un espacio métrico, entonces
  \[\{B_r(p) \colon r \in (0,\oo)\}\]
  es base en $p$.
  \[\{B_r(p) \colon r \in (0,\oo), p \in X\}\]
  es base de $\t$.
\end{example}

\begin{example}
  $X$ con la topología discreta,
  \[\{\{p\} \colon p \in X\}\]
  es base.
\end{example}

\begin{proposition}
  Sea $X \neq \varnothing$ y $\B \subseteq \P(X)$.
  $\B$ es base de alguna topología si y solo si
  \begin{enumerate}[(i)]
    \item $\B$ cubre $X$
    \item $B_1,B_2 \in \B $, $\forall p \in B_1 \cap B_2, \exists B_p \in \B$ tal que $p \in B_p \subseteq B_1 \cap B_2$.
  \end{enumerate}
\end{proposition}
\begin{proof}
  \
  \begin{itemize}
    \item \boxed{\implies}
    $(X,\t)$ es topología y $\B$ es base de $\t$.
    Como $\B$ es base para $p \in X$, entonces existe $B_p \in \B$ vecindad de $p$. Luego
    \[X \subseteq \bigcup_{p \in X} B_p \subseteq X\]
    por lo que $(i)$ se cumple, ya que
    \[X = \bigcup_{p \in X} B_p \subseteq \bigcup_{B \in \B} B \subseteq X.\]

    Veamos ahora (ii). Tomamos $B_1,B_2 \in \B$, entonces $B_1 \cap B_2$ es abierto:
    Si $B_1 \cap B_2 = \varnothing$, entonces ta está en la topología.
    Si $B_1 \cap B_2 \neq \varnothing$, entonces existe $p \in B_1 \cap B_2$. Por definición de base, existe $B$ vecindad de $p$ con $B \subseteq B_1 \cap B_2$.

    \item \boxed{\impliedby}
    Sea $\t = \{\varnothing\} \cup \{\bigcup_{\a \in I} B_\a \colon \{B_\a\}_{\a \in I} \subseteq \B\}$ el conjunto de uniones arbitrarias de elementos de $\B$. Veamos que $\t$ es topología y $\B$ es base de $\t$.
    \begin{itemize}
      \item $\varnothing \in \t$.
      \item Por (i), $X = \bigcup_{B \in \B} B \in \t$, y entonces $X \in \t$.
      \item Si tomamos $\O_\a = \bigcup_{\beta \in I_\a} B_\beta^\a \in \t$
      \[\bigcup_{\beta \in J} \O_\beta = \bigcup_{\a \in I} \bigcup_{\beta \in I_\a} B_\beta^\a \in \t.\]
      \item Sean $\O_1, \O_2 \in \t$, entonces
      \[\O_1 = \bigcup_{\a \in I_1} B_\a^1, \quad
      \O_2 = \bigcup_{\beta \in I_2} B_\b^2.\]
      Sea $p \in \O_1 \cap \O_2$, entonces $p \in \O_1$ y $p \in \O_2$.
      Luego existe $\a \in I_1$ tal que $p \in B_\a^1$ y existe $\b \in I_2$ tal que $p \in B_\b^2$. Por lo tanto $p \in B_\a^1 \cap B_\b^2$.
      Luego $\O_1 \cap \O_2 \in \t$.
    \end{itemize}
  \end{itemize}
\end{proof}

\begin{example}
  Si $(X,\t_X)$ e $(Y,\t_Y)$ son espacios topológicos, entonces
  \[\{\O_1 \times \O_2 \colon \O_1 \subseteq \t_X, \O_2 \subseteq \t_Y\}\]
  es una base para una topología en $X \times Y$.
\end{example}
\begin{proof}
  \begin{itemize}
    \item Primero notemos que $X \in \t_X$ e $Y \in \t_Y$, por lo que $X \times Y \in \B$.
    \item Tomemos $\O_1 \times S_1$ y $\O_2 \times S_2$ con $\O_i \in \t_X$, $S_i \in \t_Y$.
    \begin{align*}
      (x,y) &\in (\O_1 \times S_1) \cap (\O_2 \times S_2) \\
      &= \{(x \in \O_1 \land y \in S_1) \land (x \in \O_2 \land y \in S_2)\} \\
      &= (\O_1 \cap \O_2) \times (S_1 \cap S_2)
    \end{align*}
    \[\implies x \in \O_1 \cap \O_2 \quad\land\quad y \in S_1 \cap S_2\]
    \[\implies x \in \O_1 \land y \in S_1 \land x \in \O_2 \land y \in S_2\]
    \[\implies (x,y) \in (\O_1 \times S_1) \cup (\O_2 \times S_2)\]
    \[\implies \]
  \end{itemize}
\end{proof}
%% terminar demostración

\fecha{Clase 2. Miércoles 11 de Marzo}
\begin{definition}
  La topología dada por la base anterior se llama \wea{topología producto}.
\end{definition}

\begin{example}
  La topología producto en $\R \times \R$ es igual a la topología en $\R^2$ dada por la norma usual.
\end{example}

%% Agregar comando \mathbb{S}
%% Ver que pasa con \-

\begin{exercise}
  Sea $(X;d)$ espacio métrico. Muestre que $x \in \bar{E}$ si y solo si $\exists (x_n)_{n\in\N} \subseteq E$ tal que
  \[\lim_{n\to\oo} x_n = x.\]
\end{exercise}

\begin{definition}
  Sea $(X,\t)$ un espacio topológico. Si $C \subseteq X$, decimos que $\O$ es \wea{vecindad de $C$} si $\O\in\t$ y $C \subseteq \O$.
\end{definition}

\begin{definition}
  Sea $(X,\t)$ espacio topológico y $(a_n)_{n\in\N} \subseteq X$. Decimos que $\lim_{n\to\oo} a_n = a$ si $\forall \O$ vecindad de $a$, existe $n_0$ tal que $a_n \in \O$, $\forall n \geq n_0$.
\end{definition}

\begin{example}
  Si consideramos $X$ con la topología discreta, entonces una sucesión $(a_n)_{n\in\N}$ es convergente si y solo si $\exists n_0$ tal que $a_n = a$ para todo $n \geq n_0$. Es decir, las sucesiones que convergen son constantes desde un cierto punto en adelante.
\end{example}

\newpage
\subsection{Axiomas de Separación}

\begin{definition}
  Sea $(X,\t)$ espacio topológico. Decimos que $X$ es:
  \begin{enumerate}[(1)]
    \item \wea{Tychonoff} si $\forall u,v \in X$, existen abiertos $\O_u, \O_v \in \t$ tales que
    \[u \in \O_u \setminus \O_v, \quad v \in \O_v \setminus \O_u.\]

    \item \wea{Hausdorff} si $\forall u,v \in X$, existen $\O_u, \O_v \in \t$ tales que
    \[u \in \O_u, \quad v \in \O_v, \quad \O_u \cap \O_v = \varnothing.\]

    \item \wea{Regular} si $X$ es Tychonoff y además para $F \subseteq X$ cerrado y $u \in X \setminus F$ existen $\O, \O_u \in \t$ tales que
    \[F \subseteq \O, \quad
    u \in \O_u, \quad
    \O \cap \O_u = \varnothing.\]

    \item \wea{Normal} si $X$ es Tychonoff y si $F_1, F_2 \subseteq X$ son cerrados entonces existen $\O_1,\O_2 \in \t$ tales que
    \[F_1 \subseteq \O_1, \quad F_2 \subseteq \O_2, \quad \O_1 \cap \O_2 = \varnothing.\]
  \end{enumerate}
\end{definition}

\begin{proposition}
  $(X,\t)$ es Tychonoff si y solo si $\{v\}$ es cerrado para cualquier $v \in X$.
\end{proposition}
\begin{proof}
  \begin{itemize}
    \item \boxed{\implies} Supongamos que $(X,\t)$ es Tychonoff.
    Sea $v \in X$, y sea $u \in X \setminus \{v\}$. Luego existe $\O_u \in \t$ tal que $u \in \O_u$ y $v \notin \O_u$. Luego $u \in \O_u \subseteq X \setminus \{v\}$, y como $u$ es arbitrario, $X \setminus \{v\}$ es abierto. Luego $\{v\}$ es cerrado.

    \item \boxed{\impliedby} Supongamos que $\{v\}$ es cerrado para todo $v \in X$. Sean $u,v \in X$ tales que $u \neq v$. Entonces $u \in X \setminus \{v\}$. Como $X \setminus \{v\}$ es abierto, existe $\O_u$ vecindad de $u$ tal que $\O_u \subseteq X \setminus \{v\}$. Por lo tanto $v \notin \O_u$. De la misma forma existe $\O_v$ vecindad tal que $u \notin \O_v$ y $v \in \O_v$.
  \end{itemize}
\end{proof}

\begin{corollary}
  Normal $\subseteq$ Regular $\subseteq$ Hausdorff $\subseteq$ Tychonoff.
\end{corollary}

\begin{definition}
  Sea $(X,d)$ espacio métrico y sean $p \in X$, $C \subseteq X$. La \wea{distancia entre $p$ y $C$} es
  \[\dist(p,C) = \inf\{\dist(p,c) \colon c \in C\}.\]
\end{definition}

\begin{theorem}
  Si $(X,d)$ es un espacio métrico, entonces es normal.
\end{theorem}
\begin{proof}
  Tomemos $F_1,F_2$ cerrados y disjuntos. Definimos
  \begin{align*}
    \O_1 &= \{a \in X \colon \dist(a,F_1) < \dist(a,F_2)\}, \\
    \O_2 &= \{b \in X \colon \dist(b,F_2) < \dist(b,F_1)\}.
  \end{align*}
  Tenemos que $\O_1 \cap \O_2 = \varnothing$. Notemos que si $a \notin F_2$, entonces $\dist(a,F_2)>0$, pues si $\dist(a,F_2) = 0$ entonces existe $(a_n)_{n\in\N} \subseteq F_2$ tal que $\lim_{n\to\oo} a_n = a$, por lo que $a \in \bar{F_2} = F_2$, contradicción. Entonces $F_1 \subseteq \O_1$ y $F_2 \subseteq \O_2$.

  Veamos que $\O_i$ es abierto:
  Sea $a \in \O_1$, por demostrar que $\exists \e>0$ tal que si $d(a,b) < \e$ entonces $b \in \O_1$ (si $b \in B_\e(a)$ entonces $b \in \O$).
  Queremos ver
  \[d(b,F_1) \leq d(b,p_1), \quad \forall p_1 \in F_1.\]
  Notemos que $d(b,F_1) \leq d(a,b) + d(a,p_1)$. Luego $\forall d >0$, $\exists p_1 \in F_1$ tal que
  \[d(a,p_1) \leq \dist(a,F_1) + \d.\]
  Luego $\forall \d>0$,
  \[d(b,F_1) \leq d(a,b) + \dist(a,F_1) + \d.\]
  Tomando $\d \to 0$
  \begin{align*}
    \dist(b,F_1) &\leq d(a,b) + \dist(a,F_1) \\
    &\leq d(a,b) + \dist(a,F_1) - \dist(a,F_2) + \dist(a,F_2).
  \end{align*}
  De la misma manera:
  \[d(a,F_2) \leq d(a,b) + \dist(b,F_2)\]
  Luego
  \begin{align*}
    d(b,F_1) &\leq 2d(a,b) + \dist(a,F_1) - d(a,F_2) + \dist(b,F_2) \\
    \dist(b,F_1) - \dist(b,F_2) &\leq 2d(a,b) + \dist(a,F_1) - \dist(a,F_2).
  \end{align*}
  Tomamos $\e>0$ tal que $2\e < \dist(a,F_2) - \dist(a,F_1)$, y entonces $b \in \O$. Por lo tanto $(X,d)$ es normal.
\end{proof}

\begin{theorem}
  Sea $X$ Tychonoff. $X$ es normal si y solo si $\forall F \subseteq X$ cerrado y $U$ vecindad abierta de $F$, existe $\O$ abierto tal que
  \[F \subseteq \O \subseteq \bar\O \subseteq U.\]
\end{theorem}
\begin{proof}
  \begin{itemize}
    \item \boxed{\implies} Sea $F$ cerrado y $U$ abierto con $F \subseteq U$. Consideremos $F_2 = X \setminus U$. Tenemos que $F_2$ es cerrado y $F \cap F_2 = \varnothing$. Luego como $X$ es normal, existen $\O_1,\O_2 \in \t$ tales que
    \[F \subseteq \O_1, \quad F_2 \subseteq \O_2, \quad \O_1 \cap \O_2 = \varnothing.\]
    Veamos que$\O_1$ cumple. Ta tenemos que $F \subseteq \O_1$, además $\O_1 \cap \O_2 = \varnothing$ y $\O_2 \supseteq X \setminus U$. Por lo tanto $X \setminus \O_2 \subseteq U$. Además $\O_1 \subseteq X \setminus \O_2$
    \[\implies \bar{\O_1} \subseteq \bar{X \setminus \O_2} = X \setminus \O_2 \subseteq U.\]
    Por lo tanto $F \subseteq \O_1 \subseteq \bar{\O_1} \subseteq U$.

    \item \boxed{\impliedby} Sean $F_1, F_2$ cerrados tales que $F_1 \cap F_2 = \varnothing$. Sea $U = X \setminus F_2$, entonces $F_1 \subseteq U$.
    Por propiedad, existe $\O \in \t$ tal que
    \[F_1 \subseteq \O \subseteq \bar\O \subseteq U = X \setminus F_2.\]
    Tomamos $\O_2 = X \setminus \bar\O \in \t$. Entonces $F_2 \subseteq \O_2$. Similar para $F_1$.
  \end{itemize}
\end{proof}

\begin{definition}
  Un espacio topológico $(X,\t)$ se dice \wea{primero contable} si $\forall p \in X$ existe una base de abiertos en $p$ tal que la base es numerable.
\end{definition}

\begin{definition}
  $(X,\t)$ se dice \wea{segundo contable} si existe una base de $\t$ numerable.
\end{definition}

\begin{example}
  Todo espacio métrico $(X,d)$ es primero contable. Todo punto $p \in X$ tiene base $\{B_r(p) \colon p \in \Q_{>0}\}$.
\end{example}

\begin{definition}
  Sea $(X,\t)$ espacio topológico. Decimos que $A \subseteq X$ es \wea{denso} si $\bar{A} = X$.
\end{definition}

\begin{definition}
  Un espacio topológico $(X,\t)$ se dice \wea{separable} si existe un subconjunto denso numerable.
\end{definition}

\begin{theorem}
  Sea $(X,d)$ espacio métrico. $X$ es segundo contable si y solo si $X$ es separable.
\end{theorem}
\begin{proof}\
  \begin{itemize}
    \item \boxed{\impliedby} Tomamos $E = \{e_i\}_{i\in\N}$ denso. Tomamos $$\B = \{B_r(e_i) \colon r \in \Q_{>0}, e_i \in E\}.$$
    Veamos que es base:
    \begin{enumerate}[(i)]
      \item $\bigcup_{B\in\B} B = X$: trivial %% demostrar
      \item $B_1, B_2\in \B \implies \exists B_3 \subseteq B_1 \cap B_2$:
      
      Sean $B_1, B_2$. Tenemos que< $B_1 \cap B_2 \neq \varnothing$. Sea $p \in B_1 \cap B_2$, luego existen $r_1,r_2$ tales que $B_{r_i}(p) \subseteq B_i$. Como $E$ es denso, existe $e_i \in B_{r_i}(p)$, y como $B_{r_i}(p)$ son abiertos, existe $\e$ tal que $B_\e(e_i) \subseteq B_{r_i}(p)$. Basta con tomar $\e$ racionales.
    \end{enumerate}

    \item \boxed{\implies} Sea $\{B_i\}_{i\in\N}$ base de abiertos numerable. Basta tomar $e_i \in B_i$, y considerar $E = \{e_i\}_{i\in\N}$.
  \end{itemize}
\end{proof}

\begin{theorem}
  Si $(X,\t)$ es un espacio topológico primero contable, entonces $a \in \bar{E}$ si y solo si $\exists (a_n)_{n\in\N} \subseteq E$ tal que $\lim_{n\to\oo} a_n = a$.
\end{theorem}
\begin{proof}\
  \begin{itemize}
    \item \boxed{\implies} Sea $p \in \bar{E}$. Tomamos $\{B_n\}$ base de abiertos de $p$ numerable. Notemos que, como $p \in \bar{E}$, entonces $B_1 \cap E \neq \varnothing$. Tomamos $p_1 \in B_1 \cap E$. Luego
    \[p_i \in E \cap \bigcap_{j\leq i} B_i \neq \varnothing.\]
    Este conjunto es no vacío porque es abierto y $p \in B_1 \cap \dots \cap B_i$.
    Ahora veamos que $p_n \to p$.

    Tomamos $B$ abierto con $p \in B$. Como $\{B_i\}$ es base, existe $i_0$ tal que $B_{i_0} \subseteq B$. Luego para todo $i \geq i_0$
    \[p_i \in B_i \cap \dots \cap B_{i_0} \cap \dots \cap B_1 \cap E.\]
    Luego $p_i \in B_{i_0}$. Para todo $i \geq i_0$, $p_i \in B$.
  \end{itemize}
\end{proof}

\begin{theorem}
  [Urysohn]
  Si $(X,\t)$ es normal y segundo contable, entonces $(X,\t)$ es metrizable.
\end{theorem}

%\fecha{Clase 3.}
%\fecha{Clase 4.}
%\fecha{Clase 5.}
%\fecha{Clase 6.}
%\fecha{Clase 7.}
%\fecha{Clase 8.}
%\fecha{Clase 9.}
%\fecha{Clase 10.}
%\fecha{Clase 11.}
%\fecha{Clase 12.}
%\fecha{Clase 13.}
%\fecha{Clase 14.}
%\fecha{Clase 15.}
%\fecha{Clase 16.}
%\fecha{Clase 17.}
%\fecha{Clase 18.}
%\fecha{Clase 19.}
%\fecha{Clase 20.}
%\fecha{Clase 21.}
%\fecha{Clase 22.}
%\fecha{Clase 23.}
%\fecha{Clase 24.}
%\fecha{Clase 25.}
%\fecha{Clase 26.}
%\fecha{Clase 27.}
%\fecha{Clase 28.}
%\fecha{Clase 29.}
%\fecha{Clase 30.}

\end{document}